\documentclass[11pt,aspectratio=169]{beamer}
\usetheme{SimplePlus}

\usepackage[T1]{fontenc}
\usepackage[utf8]{inputenc}
\usepackage{amsmath,amssymb,mathtools}
\usepackage{booktabs}
\usepackage{array}
\usepackage{appendixnumberbeamer}
\usepackage{hyperref}

\newcommand{\E}{\mathbb{E}}
\newcommand{\Var}{\mathrm{Var}}
\newcommand{\Cov}{\mathrm{Cov}}
\newcommand{\Prob}{\mathbb{P}}
\newcommand{\ATT}{\mathrm{ATT}}
\newcommand{\ATE}{\mathrm{ATE}}
\newcommand{\indep}{\perp\!\!\!\perp}
\newcommand{\dd}{\mathrm{d}}
\newcommand{\derivbutton}[1]{\vspace{0.4em}\hfill\hyperlink{app:#1}{\beamerbutton{Appendix details}}}
\newcommand{\backbutton}[1]{\vspace{0.4em}\hfill\hyperlink{main:#1}{\beamerbutton{Back to main slide}}}

\title{Applied Econometrics: Session 7}
\subtitle{Difference-in-Differences, Panel Data, and Fixed Effects}
\author{PhD Intensive Course}
\institute{Lecture 7}
\date{\today}

\begin{document}

\begin{frame}
  \titlepage
\end{frame}

\begin{frame}{How This Lecture Fits the Course}
  Earlier lectures asked how to make treated and untreated units comparable.
  \vspace{0.4em}

  Today we use a different source of comparison:
  \begin{block}{Core idea}
    If treated and untreated units would have followed the same trend without treatment, then the control group's change can proxy for the treated group's missing counterfactual change.
  \end{block}

  \begin{itemize}
    \item Difference-in-differences uses comparisons across groups and across time.
    \item Panel fixed effects use repeated observations to remove fixed unit differences.
    \item Both methods are powerful, but only under clear trend assumptions.
  \end{itemize}
\end{frame}

\begin{frame}{Roadmap for Today}
  \tableofcontents
\end{frame}

\section{The 2-by-2 DiD Design}

\begin{frame}{Motivating Example: A City Policy}
  Suppose one city introduces a new job-search program in 2025.
  We observe average employment in:
  \begin{itemize}
    \item the treated city before and after the program;
    \item a similar control city before and after the program.
  \end{itemize}
  \vspace{0.3em}
  \begin{block}{Causal question}
    What would employment in the treated city have been after 2025 if the program had not happened?
  \end{block}
  \begin{block}{Why this is hard}
    We observe the treated city after treatment, but not the same treated city after no treatment.
  \end{block}
\end{frame}

\begin{frame}{The Four Cells}
  Let \(D_i=1\) denote the treated group and \(Post_t=1\) denote the period after treatment begins.
  \[
  \begin{array}{c c c}
  \toprule
  & Post_t=0 & Post_t=1 \\
  \midrule
  D_i=0\ \text{control} & \bar Y_{00} & \bar Y_{01} \\
  D_i=1\ \text{treated} & \bar Y_{10} & \bar Y_{11} \\
  \bottomrule
  \end{array}
  \]
  \vspace{0.6em}
  \begin{itemize}
    \item \(\bar Y_{10}\): treated group before treatment.
    \item \(\bar Y_{11}\): treated group after treatment.
    \item \(\bar Y_{00},\bar Y_{01}\): control group before and after.
  \end{itemize}
\end{frame}

\begin{frame}{Two Naive Comparisons}
  \textbf{Before-after for treated:}
  \[
    \bar Y_{11}-\bar Y_{10}.
  \]
  This is biased if the outcome would have changed over time even without treatment.

  \vspace{0.8em}
  \textbf{Post-period treated-control gap:}
  \[
    \bar Y_{11}-\bar Y_{01}.
  \]
  This is biased if treated and control groups differ in baseline levels.

  \begin{block}{DiD's move}
    Use the control group's change over time to remove common trends, and use the treated group's pre-period level to remove fixed baseline differences.
  \end{block}
\end{frame}

\begin{frame}{The Difference-in-Differences Estimator}
  \hypertarget{main:did-estimator}{}
  The 2-by-2 DiD estimator is
  \[
    \hat\delta_{\mathrm{DiD}}
    =
    (\bar Y_{11}-\bar Y_{10})
    -
    (\bar Y_{01}-\bar Y_{00}).
  \]
  Equivalent form:
  \[
    \hat\delta_{\mathrm{DiD}}
    =
    (\bar Y_{11}-\bar Y_{01})
    -
    (\bar Y_{10}-\bar Y_{00}).
  \]
  \begin{itemize}
    \item First expression: difference in time changes.
    \item Second expression: difference in group gaps.
    \item Same number, two useful interpretations.
  \end{itemize}
  \derivbutton{did-algebra}
\end{frame}

\begin{frame}{Numeric Example}
  \[
  \begin{array}{lcc}
  \toprule
  & \text{Before} & \text{After} \\
  \midrule
  \text{Control city} & 50 & 58 \\
  \text{Treated city} & 40 & 55 \\
  \bottomrule
  \end{array}
  \]
  \vspace{0.4em}

  \[
    \text{Treated change}=55-40=15,
    \qquad
    \text{Control change}=58-50=8.
  \]
  \[
    \hat\delta_{\mathrm{DiD}}=15-8=7.
  \]
  \begin{block}{Interpretation}
    Employment rose by 15 in the treated city, but 8 of those points look like a common trend. The remaining 7 is the DiD estimate.
  \end{block}
\end{frame}

\begin{frame}{Counterfactual Picture in Words}
  DiD constructs the missing treated-after-without-treatment outcome:
  \[
    \widehat{\E[Y_{0}\mid D=1,Post=1]}
    =
    \bar Y_{10}
    +
    (\bar Y_{01}-\bar Y_{00}).
  \]
  \vspace{0.4em}
  In the numeric example:
  \[
    \widehat{Y}_{\text{treated, after, no treatment}}
    =
    40+(58-50)=48.
  \]
  We actually observe \(55\), so the estimated effect is
  \[
    55-48=7.
  \]
  \begin{block}{Student translation}
    Take the treated group's starting point. Add the control group's growth. Compare the observed treated-after outcome to that counterfactual.
  \end{block}
\end{frame}

\begin{frame}{Potential Outcomes Notation}
  \hypertarget{main:did-po}{}
  Let \(Y_{it}(1)\) be unit \(i\)'s outcome at time \(t\) if treated, and \(Y_{it}(0)\) if untreated.
  Treatment starts only in the post period for the treated group.

  \[
    Y_{it}=Y_{it}(0)+D_i Post_t\left(Y_{it}(1)-Y_{it}(0)\right).
  \]

  The target is the average treatment effect on the treated after treatment:
  \[
    \ATT
    =
    \E[Y_{i1}(1)-Y_{i1}(0)\mid D_i=1].
  \]

  The missing object is
  \[
    \E[Y_{i1}(0)\mid D_i=1].
  \]
  \derivbutton{did-identification}
\end{frame}

\begin{frame}{Parallel Trends}
  \hypertarget{main:parallel-trends}{}
  Parallel trends says that, absent treatment, treated and control groups would have had the same average change:
  \[
    \E[Y_{i1}(0)-Y_{i0}(0)\mid D_i=1]
    =
    \E[Y_{i1}(0)-Y_{i0}(0)\mid D_i=0].
  \]
  \vspace{0.4em}
  This does \textbf{not} require:
  \begin{itemize}
    \item equal levels before treatment;
    \item equal covariate distributions;
    \item random assignment of treatment.
  \end{itemize}
  \vspace{0.2em}
  It requires equal untreated \textbf{trends}.
  \derivbutton{did-identification}
\end{frame}

\begin{frame}{What Parallel Trends Allows}
  Parallel trends allows fixed differences in levels:
  \[
    \E[Y_{i0}(0)\mid D=1]
    \neq
    \E[Y_{i0}(0)\mid D=0].
  \]

  \vspace{0.5em}
  Example:
  \begin{itemize}
    \item Rich cities can have higher employment than poor cities.
    \item That is not fatal if both would have grown by the same amount without treatment.
    \item DiD subtracts out the level difference by comparing changes.
  \end{itemize}
  \begin{block}{Important distinction}
    A control group can be useful even if it starts above or below the treated group. It becomes dangerous when it is on a different path.
  \end{block}
\end{frame}

\begin{frame}{What Breaks Parallel Trends}
  Parallel trends fails when treatment timing is related to untreated growth.
  \vspace{0.3em}
  Common examples:
  \begin{itemize}
    \item A city receives a job program because unemployment was already worsening.
    \item A firm advertises in markets where sales were already falling.
    \item A school adopts a curriculum because test scores had begun rising.
    \item A policy is timed to a local shock that affects the outcome.
  \end{itemize}
  \begin{block}{Diagnostic question}
    Would the treated group have changed like the control group if the treatment had not happened?
  \end{block}
\end{frame}

\begin{frame}{Regression Version of 2-by-2 DiD}
  \hypertarget{main:did-reg}{}
  Estimate
  \[
    Y_{it}
    =
    \alpha
    +
    \gamma D_i
    +
    \lambda Post_t
    +
    \delta (D_i\times Post_t)
    +
    u_{it}.
  \]
  \vspace{0.3em}
  The coefficient \(\delta\) equals the 2-by-2 DiD estimator.
  \[
    \hat\delta
    =
    (\bar Y_{11}-\bar Y_{10})
    -
    (\bar Y_{01}-\bar Y_{00}).
  \]
  \begin{itemize}
    \item \(D_i\): permanent treated-control level difference.
    \item \(Post_t\): common post-period shock.
    \item \(D_i\times Post_t\): treatment turns on for treated units after policy.
  \end{itemize}
  \derivbutton{did-regression}
\end{frame}

\begin{frame}{Why Regression Is Convenient}
  The regression form gives the same point estimate in the simple 2-by-2 design, but adds:
  \begin{itemize}
    \item standard errors;
    \item covariates, if pre-treatment and justified;
    \item many periods and many units;
    \item fixed effects;
    \item event-study coefficients.
  \end{itemize}
  \begin{block}{Warning}
    Regression does not make parallel trends true. It only gives a convenient way to estimate the design.
  \end{block}
\end{frame}

\section{Event-Study Intuition}

\begin{frame}{Why More Pre-Periods Help}
  With only one pre-period, parallel trends is not visually testable.
  With several pre-periods, we can ask:
  \[
    \text{Were treated and control groups already moving differently before treatment?}
  \]
  \vspace{0.3em}
  This is not a proof of parallel trends, because the unobserved post-treatment counterfactual is still missing.
  \vspace{0.3em}
  But it is an important credibility check.
  \begin{block}{Practical rule}
    If treated and control groups diverge before treatment, be cautious about interpreting the post-treatment divergence as causal.
  \end{block}
\end{frame}

\begin{frame}{Event Time}
  Event studies re-index time relative to treatment:
  \[
    k = t - T_i^\ast,
  \]
  where \(T_i^\ast\) is the treatment start date.
  \vspace{0.4em}
  \begin{itemize}
    \item \(k=-3\): three periods before treatment.
    \item \(k=-1\): the omitted reference period, usually just before treatment.
    \item \(k=0\): first treated period.
    \item \(k=2\): two periods after treatment.
  \end{itemize}
  \begin{block}{Goal}
    Estimate how the treated-control difference evolves before and after the treatment date.
  \end{block}
\end{frame}

\begin{frame}{A Simple Event-Study Regression}
  \hypertarget{main:event-study}{}
  With one treated cohort and never-treated controls:
  \[
    Y_{it}
    =
    \alpha_i+\lambda_t
    +
    \sum_{k\neq -1}\delta_k
    \left(1\{t-T^\ast_i=k\}\times D_i\right)
    +
    u_{it}.
  \]
  \vspace{0.3em}
  \begin{itemize}
    \item Unit fixed effects absorb permanent treated-control differences.
    \item Time fixed effects absorb common shocks.
    \item \(\delta_k\) compares treated units to controls at event time \(k\), relative to \(k=-1\).
  \end{itemize}
  \derivbutton{event-study}
\end{frame}

\begin{frame}{How to Read Event-Study Coefficients}
  \begin{itemize}
    \item Pre-treatment coefficients (\(k<0\)) are called leads.
    \item Post-treatment coefficients (\(k\geq 0\)) are called lags or dynamic effects.
    \item Leads near zero support the idea that groups were not already diverging.
    \item A jump at \(k=0\) suggests an immediate treatment effect.
    \item A gradual increase after \(k=0\) suggests delayed adjustment.
  \end{itemize}
  \begin{block}{Student warning}
    Do not say "the pre-trend test proves parallel trends." It can only fail to detect obvious pre-treatment violations.
  \end{block}
\end{frame}

\begin{frame}{Event-Study Pitfalls}
  \begin{itemize}
    \item Few pre-periods mean weak diagnostics.
    \item Wide confidence intervals can hide economically important pre-trends.
    \item Anticipation effects can move outcomes before treatment begins.
    \item Changing composition can make the treated group different over time.
    \item With staggered adoption and heterogeneous effects, simple TWFE event studies can be misleading.
  \end{itemize}
  \begin{block}{For this course}
    Today we focus on the clean case: one treated cohort, never-treated controls, and a common treatment start date.
  \end{block}
\end{frame}

\section{Panel Data}

\begin{frame}{What Is Panel Data?}
  Panel data repeatedly observes the same unit over time:
  \[
    i=1,\ldots,N,
    \qquad
    t=1,\ldots,T.
  \]
  Examples:
  \begin{itemize}
    \item people observed in multiple survey waves;
    \item firms observed across years;
    \item municipalities observed monthly;
    \item schools observed before and after a reform.
  \end{itemize}
  \begin{block}{Why panels matter}
    Repeated observations let us compare a unit to itself over time, not only to other units.
  \end{block}
\end{frame}

\begin{frame}{Panel Notation}
  \[
    Y_{it} = \text{outcome for unit } i \text{ in period } t.
  \]
  \[
    D_{it} = \text{treatment status for unit } i \text{ in period } t.
  \]
  \[
    X_{it} = \text{observed covariates for unit } i \text{ in period } t.
  \]

  \vspace{0.4em}
  Panel regressions must respect two dimensions:
  \begin{itemize}
    \item observations from the same unit are not independent over time;
    \item shocks in the same period can affect many units at once.
  \end{itemize}
\end{frame}

\begin{frame}{Why Panels Help with Unobserved Confounding}
  Suppose ability, geography, management quality, or culture affects the outcome.
  If that feature is fixed for unit \(i\), write it as \(a_i\).
  \[
    Y_{it} = \beta D_{it}+a_i+u_{it}.
  \]
  \vspace{0.3em}
  \begin{itemize}
    \item Cross-sectional OLS confuses \(D_{it}\) with \(a_i\) if treated units are permanently different.
    \item Fixed effects compare each unit to its own average.
    \item Anything constant within the unit is removed.
  \end{itemize}
  \begin{block}{Plain English}
    Fixed effects control for stable differences that we may not observe but that travel with the unit over time.
  \end{block}
\end{frame}

\begin{frame}{The Unit Fixed-Effects Model}
  \hypertarget{main:unit-fe}{}
  A unit fixed-effects model is
  \[
    Y_{it}=\beta D_{it}+a_i+u_{it}.
  \]
  \(a_i\) is a unit-specific intercept.
  \vspace{0.3em}
  \begin{itemize}
    \item Each person, firm, school, or city gets its own baseline.
    \item Identification comes from within-unit changes in treatment.
    \item Units whose treatment never changes do not identify \(\beta\), though they may help estimate time effects in richer models.
  \end{itemize}
  \derivbutton{within-oneway}
\end{frame}

\begin{frame}{The Within Transformation}
  \hypertarget{main:within}{}
  For each unit define time averages:
  \[
    \bar Y_i=\frac{1}{T_i}\sum_t Y_{it},
    \qquad
    \bar D_i=\frac{1}{T_i}\sum_t D_{it}.
  \]
  Subtract the unit average:
  \[
    Y_{it}-\bar Y_i
    =
    \beta(D_{it}-\bar D_i)
    +
    (u_{it}-\bar u_i).
  \]
  The fixed unit term disappears because \(a_i-\bar a_i=0\).
  \begin{block}{Mechanical meaning}
    FE regression is OLS on demeaned outcomes and demeaned treatment.
  \end{block}
  \derivbutton{within-oneway}
\end{frame}

\begin{frame}{What Fixed Effects Cannot Estimate}
  In a unit FE model, time-invariant variables are absorbed:
  \[
    Z_i-\bar Z_i=0.
  \]
  Examples:
  \begin{itemize}
    \item sex, birth cohort, ethnicity, country of birth;
    \item distance to the coast, historical institutions, baseline geography;
    \item any variable measured once and copied across all years.
  \end{itemize}
  \begin{block}{Interpretation}
    FE can control for time-invariant variables, but cannot estimate their separate coefficients in the same model.
  \end{block}
\end{frame}

\begin{frame}{Time Fixed Effects}
  A time fixed effect is a period-specific intercept:
  \[
    Y_{it}=\beta D_{it}+a_i+\lambda_t+u_{it}.
  \]
  \vspace{0.3em}
  \(\lambda_t\) absorbs shocks common to all units in period \(t\):
  \begin{itemize}
    \item national recessions;
    \item inflation;
    \item pandemics;
    \item common seasonality;
    \item changes in national measurement rules.
  \end{itemize}
  \begin{block}{Two-way fixed effects}
    Unit FE compare units to themselves; time FE remove shocks shared by everyone at a point in time.
  \end{block}
\end{frame}

\begin{frame}{Two-Way Fixed Effects and DiD}
  \hypertarget{main:twfe}{}
  The basic TWFE DiD regression is
  \[
    Y_{it}
    =
    a_i+\lambda_t+\delta D_{it}+u_{it}.
  \]
  In a simple 2-by-2 design:
  \[
    D_{it}=1\{i\text{ treated}\}\times 1\{t\text{ post}\}.
  \]
  Then \(\hat\delta\) equals the 2-by-2 DiD estimate.
  \begin{itemize}
    \item \(a_i\) replaces the treated-group dummy.
    \item \(\lambda_t\) replaces the post dummy.
    \item \(D_{it}\) is the interaction that turns treatment on.
  \end{itemize}
  \derivbutton{twfe-equivalence}
\end{frame}

\begin{frame}{Simple TWFE Equivalence}
  Compare the two regressions:
  \[
    Y_{it}
    =
    \alpha+\gamma D_i+\lambda Post_t+\delta(D_iPost_t)+u_{it}
  \]
  and
  \[
    Y_{it}
    =
    a_i+\lambda_t+\delta D_{it}+u_{it}.
  \]
  In the 2-by-2 case, \(D_{it}=D_iPost_t\).
  \vspace{0.3em}
  \begin{block}{Same estimator, different notation}
    The first expression is easier for beginners. The second scales to many units and many periods.
  \end{block}
\end{frame}

\begin{frame}{Fixed-Effect Demeaning with Two-Way FE}
  For balanced panels, two-way demeaning removes unit and time averages:
  \[
    \tilde Y_{it}
    =
    Y_{it}-\bar Y_i-\bar Y_t+\bar Y.
  \]
  \[
    \tilde D_{it}
    =
    D_{it}-\bar D_i-\bar D_t+\bar D.
  \]
  Then TWFE estimates \(\delta\) by regressing \(\tilde Y_{it}\) on \(\tilde D_{it}\).
  \begin{block}{Intuition}
    TWFE uses only the outcome variation left after removing each unit's average and each time period's average.
  \end{block}
\end{frame}

\section{Practical Diagnostics}

\begin{frame}{Checklist Before Running DiD}
  \begin{enumerate}
    \item Define the treatment date and treatment group precisely.
    \item Plot raw means for treated and control groups over time.
    \item Check whether pre-treatment trends look similar.
    \item Explain why treatment timing was not driven by unobserved outcome trends.
    \item Decide the level at which treatment varies and cluster standard errors there.
  \end{enumerate}
  \begin{block}{Design first}
    A table of coefficients is not a substitute for a credible comparison group.
  \end{block}
\end{frame}

\begin{frame}{What to Plot}
  Useful figures:
  \begin{itemize}
    \item treated and control mean outcomes over time;
    \item the treated-control gap over time;
    \item event-study coefficients with confidence intervals;
    \item number of units observed by group and period;
    \item distribution of treatment timing if adoption is staggered.
  \end{itemize}
  \begin{block}{What you are looking for}
    Is there a visible break at treatment time, or were groups already moving apart before treatment?
  \end{block}
\end{frame}

\begin{frame}{Standard Errors}
  Repeated observations create dependence:
  \[
    u_{it}\ \text{and}\ u_{is}
    \quad \text{may be correlated for the same unit } i.
  \]
  \vspace{0.3em}
  Practical guidance:
  \begin{itemize}
    \item cluster at the treatment-assignment level when possible;
    \item if policy varies by state, cluster by state;
    \item if policy varies by school, cluster by school;
    \item with very few clusters, conventional cluster-robust standard errors can be unreliable.
  \end{itemize}
  \begin{block}{Core point}
    DiD standard errors must respect the level at which shocks and treatment assignment are correlated.
  \end{block}
\end{frame}

\begin{frame}{Covariates in DiD}
  Covariates can help precision or adjust for composition, but they do not rescue a bad design.
  \vspace{0.3em}
  Safer covariates:
  \begin{itemize}
    \item pre-treatment characteristics;
    \item time-varying controls that are clearly not affected by treatment.
  \end{itemize}
  Risky covariates:
  \begin{itemize}
    \item post-treatment outcomes or mediators;
    \item variables that treatment itself changes;
    \item controls chosen after looking at the desired coefficient.
  \end{itemize}
\end{frame}

\begin{frame}{When Panel Fixed Effects Do Not Solve the Problem}
  Fixed effects remove only time-invariant unobserved confounding.
  They do not remove:
  \begin{itemize}
    \item unobserved shocks that change over time within units;
    \item reverse causality;
    \item anticipation before treatment;
    \item spillovers from treated to control units;
    \item measurement changes that differ by group;
    \item differential trends not captured by the model.
  \end{itemize}
  \begin{block}{Honest conclusion}
    Panel data is useful because it changes the identifying assumption. It does not eliminate the need for an identifying assumption.
  \end{block}
\end{frame}

\begin{frame}{Common Reporting Template}
  A clear DiD write-up should state:
  \begin{itemize}
    \item outcome, treatment, treated group, control group, and time window;
    \item the exact regression specification;
    \item the identifying assumption in words;
    \item evidence on pre-trends;
    \item standard-error clustering level;
    \item robustness checks and limitations.
  \end{itemize}
  \begin{block}{One-sentence identification claim}
    "Absent the policy, treated and control units would have experienced the same average change in the outcome over the study window."
  \end{block}
\end{frame}

\begin{frame}{Python Example for This Lecture}
  The companion script is:
  \[
    \texttt{lectures/code/lecture-7-did-panel.py}
  \]
  It simulates:
  \begin{itemize}
    \item a panel where parallel trends holds;
    \item a panel where treated units have a different untreated trend;
    \item 2-by-2 DiD estimates;
    \item TWFE estimates;
    \item event-study coefficients.
  \end{itemize}
  \begin{block}{What students should notice}
    The estimator works in the first simulation and is biased in the second even though the regression runs normally in both.
  \end{block}
\end{frame}

\section{Synthesis}

\begin{frame}{Main Takeaways}
  \begin{enumerate}
    \item DiD estimates treatment effects by comparing changes, not levels.
    \item The key identifying assumption is parallel untreated trends.
    \item Event studies make the trend assumption visible, but do not prove it.
    \item Unit fixed effects remove stable unit differences.
    \item Time fixed effects remove common period shocks.
    \item In the clean 2-by-2 case, TWFE and the DiD formula give the same estimate.
  \end{enumerate}
\end{frame}

\begin{frame}{Source and Further Reading}
  This lecture is based on ideas from:
  \begin{itemize}
    \item Matheus Facure Alves, \emph{Causal Inference for the Brave and True}, Chapter 13, "Difference-in-Differences".
    \item Matheus Facure Alves, \emph{Causal Inference for the Brave and True}, Chapter 14, "Panel Data and Fixed Effects".
    \item Angrist and Pischke, \emph{Mostly Harmless Econometrics}, chapters on panel data and differences-in-differences.
    \item Cunningham, \emph{Causal Inference: The Mixtape}, chapters on DiD and fixed effects.
  \end{itemize}
  \vspace{0.3em}
  Handbook links:
  \begin{itemize}
    \item \url{https://matheusfacure.github.io/python-causality-handbook/13-Difference-in-Differences.html}
    \item \url{https://matheusfacure.github.io/python-causality-handbook/14-Panel-Data-and-Fixed-Effects.html}
  \end{itemize}
\end{frame}

\begin{frame}{What Students Should Be Able to Do}
  After this lecture, students should be able to:
  \begin{itemize}
    \item compute a 2-by-2 DiD estimate by hand;
    \item state the parallel trends assumption in words and symbols;
    \item explain the purpose of event-study lead coefficients;
    \item write a panel regression with unit and time fixed effects;
    \item explain why demeaning removes unit fixed effects;
    \item describe practical diagnostics before trusting a DiD estimate.
  \end{itemize}
\end{frame}

\appendix

\begin{frame}{Appendix Map}
  \begin{itemize}
    \item A1: 2-by-2 DiD algebra.
    \item A2: DiD identification under parallel trends.
    \item A3: Regression coefficient equals DiD.
    \item A4: Event-study coefficients.
    \item A5: Within transformation for one-way fixed effects.
    \item A6: TWFE equivalence in the 2-by-2 design.
  \end{itemize}
\end{frame}

\begin{frame}{Appendix A1: 2-by-2 DiD Algebra}
  \hypertarget{app:did-algebra}{}
  Start from the change comparison:
  \[
    (\bar Y_{11}-\bar Y_{10})
    -
    (\bar Y_{01}-\bar Y_{00}).
  \]
  Rearrange terms:
  \[
    \bar Y_{11}-\bar Y_{10}-\bar Y_{01}+\bar Y_{00}
    =
    \bar Y_{11}-\bar Y_{01}-\bar Y_{10}+\bar Y_{00}.
  \]
  Group as post gap minus pre gap:
  \[
    (\bar Y_{11}-\bar Y_{01})
    -
    (\bar Y_{10}-\bar Y_{00}).
  \]
  The two formulas are numerically identical.
  \backbutton{did-estimator}
\end{frame}

\begin{frame}{Appendix A2: DiD Identification Under Parallel Trends}
  \hypertarget{app:did-identification}{}
  The target is
  \[
    \ATT
    =
    \E[Y_{i1}(1)-Y_{i1}(0)\mid D_i=1].
  \]
  The observed treated change is
  \[
  \begin{aligned}
    \E[Y_{i1}\mid D=1]-\E[Y_{i0}\mid D=1]
    &=
    \E[Y_{i1}(1)\mid D=1]-\E[Y_{i0}(0)\mid D=1] \\
    &=
    \ATT
    +
    \E[Y_{i1}(0)-Y_{i0}(0)\mid D=1].
  \end{aligned}
  \]
  The treated change equals the causal effect plus the untreated trend for treated units.
\end{frame}

\begin{frame}{Appendix A2 Continued: Subtract the Control Trend}
  Under parallel trends,
  \[
    \E[Y_{i1}(0)-Y_{i0}(0)\mid D=1]
    =
    \E[Y_{i1}(0)-Y_{i0}(0)\mid D=0].
  \]
  Since controls are untreated in both periods,
  \[
    \E[Y_{i1}(0)-Y_{i0}(0)\mid D=0]
    =
    \E[Y_{i1}-Y_{i0}\mid D=0].
  \]
  Therefore
  \[
    \ATT
    =
    \left\{\E[Y_{i1}\mid D=1]-\E[Y_{i0}\mid D=1]\right\}
    -
    \left\{\E[Y_{i1}\mid D=0]-\E[Y_{i0}\mid D=0]\right\}.
  \]
  \backbutton{parallel-trends}
\end{frame}

\begin{frame}{Appendix A3: Regression Coefficient Equals DiD}
  \hypertarget{app:did-regression}{}
  In the saturated 2-by-2 regression
  \[
    Y_{it}
    =
    \alpha+\gamma D_i+\lambda Post_t+\delta(D_iPost_t)+u_{it},
  \]
  predicted cell means are:
  \[
  \begin{array}{c c c}
  \toprule
  & Post=0 & Post=1 \\
  \midrule
  D=0 & \alpha & \alpha+\lambda \\
  D=1 & \alpha+\gamma & \alpha+\gamma+\lambda+\delta \\
  \bottomrule
  \end{array}
  \]
  \vspace{0.3em}
  The treated change is \(\lambda+\delta\). The control change is \(\lambda\).
  Difference:
  \[
    (\lambda+\delta)-\lambda=\delta.
  \]
  \backbutton{did-reg}
\end{frame}

\begin{frame}{Appendix A4: Event-Study Coefficients}
  \hypertarget{app:event-study}{}
  Let \(k=t-T^\ast_i\) be event time and omit \(k=-1\).
  \[
    Y_{it}
    =
    \alpha_i+\lambda_t
    +
    \sum_{k\neq -1}\delta_k
    1\{k=t-T^\ast_i\}D_i
    +
    u_{it}.
  \]
  The omitted period \(k=-1\) is the reference. Each \(\delta_k\) is interpreted relative to that period:
  \[
    \delta_k
    =
    \left(\text{treated-control gap at event time } k\right)
    -
    \left(\text{treated-control gap at } k=-1\right),
  \]
  after removing unit and calendar-time fixed effects.
  \backbutton{event-study}
\end{frame}

\begin{frame}{Appendix A5: Within Transformation}
  \hypertarget{app:within-oneway}{}
  Start with
  \[
    Y_{it}=\beta D_{it}+a_i+u_{it}.
  \]
  Average over time within unit \(i\):
  \[
    \bar Y_i=\beta \bar D_i+a_i+\bar u_i.
  \]
  Subtract:
  \[
  \begin{aligned}
    Y_{it}-\bar Y_i
    &=
    \beta(D_{it}-\bar D_i)+(a_i-a_i)+(u_{it}-\bar u_i)\\
    &=
    \beta(D_{it}-\bar D_i)+(u_{it}-\bar u_i).
  \end{aligned}
  \]
  The unit fixed effect \(a_i\) is gone.
  \backbutton{within}
\end{frame}

\begin{frame}{Appendix A6: TWFE Equivalence in the 2-by-2 Design}
  \hypertarget{app:twfe-equivalence}{}
  With two groups and two periods, write
  \[
    D_{it}=G_iPost_t,
  \]
  where \(G_i=1\) for treated units.
  The TWFE model is
  \[
    Y_{it}=a_i+\lambda_t+\delta D_{it}+u_{it}.
  \]
  If we replace unit fixed effects by a group dummy and time fixed effects by a post dummy, this becomes
  \[
    Y_{it}=\alpha+\gamma G_i+\lambda Post_t+\delta(G_iPost_t)+u_{it}.
  \]
  By Appendix A3, the coefficient on \(G_iPost_t\) is exactly the DiD estimator.
  \backbutton{twfe}
\end{frame}

\end{document}
